\clearpage
\section*{September 22, 2026: Which source supplies each filing's units}
\textit{Agent-drafted research entry.}

Until this correction, historical proposals filed in 2019--2022 were selected
by their initial DOB filing date but took their proposed Class~A units, lots
and building identifiers from Housing Database release 25Q4, published after
485-x took effect. A design amended after adoption therefore entered the
pre-policy distribution with its later unit count. The clearest case is 370
Livingston Street, job \texttt{B00781447}, first filed on September 28, 2022.
The official 23Q4 archive records it at 144 units; 25Q4 records 99. A recorded
declaration, document 2026010400010003, physical page~16, reproduces an April
2025 HPD application for the same job with 99 units and intended 485-x
Option~B treatment. The 99 in the historical panel was a post-policy design,
not a pre-policy choice. These sources do not date the amendment or establish
why it was made.

The problem was not limited to Livingston. A diagnostic that held panel
membership fixed and compared each historical constituent with the 23Q4 files
matched 1,902 of 1,904 constituents; 225 had different unit counts, affecting
218 parents. In the then 554-parent historical weighting sample, 109 of 592
constituents differed, affecting 106 parents. Four of the six historical
filings at exactly 99 units had other counts in 23Q4: Livingston 144,
Lafayette 104, East 198th Street 64 and Vyse Avenue 54. The other two were 99
in both releases. Differences between releases may reflect amendments or
administrative corrections; they are not estimates of a policy response. They
showed that the historical comparison needed a single dated source recorded
before adoption.

\paragraph{The unit rule now.}
Unit counts come from administrative sources according to filing period:

\begin{center}
\small
\begin{tabular}{p{0.22\textwidth}p{0.70\textwidth}}
\hline
Filings & Unit source \\
\hline
Historical, 2019--2022 &
Inactive-inclusive 23Q4 Housing Database archive: Class~A units, filing date,
lot, building identifier (BIN), description and ownership category. Archived
status and membership in the active/completed file are kept as separate
fields. A filing missing from 23Q4 is not filled from a later release. \\
Post-policy, January 1, 2025--July 8, 2026 &
25Q4 Housing Database Class~A units; DOB initial-filing units only when the
Housing Database count is missing. \\
Both periods &
A recorded zero stays zero. Documentary unit schedules and the July 2026 DOB
counts are comparison fields and never override the selected count. \\
\hline
\end{tabular}
\end{center}

Jacob selected the inactive-inclusive archive, which keeps every recorded
proposal together with its status. The archive's latest DOB update is January
12, 2024. It therefore records designs as they stood before 485-x adoption,
including any amendments made after the initial filing.

The zero rule came from a separate error found during site review. Metropolitan
Park, \texttt{Q01274555}, had entered both residential panels and the
reweighting sample with 1,000 units. DCP's 25Q4 file records zero Class~A units
and 1,000 hotel rooms, and the state-filed casino application confirms the
hotel count. Cohort construction had filtered Housing Database records on the
six-unit minimum and land-data availability before choosing a unit source, so
the recorded zero became a missing match and the positive DOB count was used
instead. The initial DOB description said ``gaming facility,'' so the separate
hotel keyword check did not catch it. The producer now reads Class~A units
from the unfiltered staged Housing Database table, and a final check requires
the selected units to equal that priority count. Zero-Class-A filings remain
in membership with their DOB counts, links and filing and refiling dates, but
contribute no residential units. The rule does not treat community-facility
beds as apartments or substitute an outside documentary count.

That correction removed four entirely zero-Class-A parents from the
post-period sample: Metropolitan Park, Mermaid, Meserole and Jamaica. On
September 20, before the historical source change, the 50-plus-unit
post-period reweighting sample fell from 316 to 312 parents and from 43,807
to 42,341 units, and parents above 300 units fell from 25 to 24. One Bronx parent became 69 rather than 69+27 units,
and a separate filing changed from 81 to 80. The effect was in the upper tail
and in the calibration composition. At that time the post-period sample still
had 39 parents at exactly 99 and nine 99+99 pairs.

\paragraph{Withdrawn alternatives and reviewed filing roles.}
Keeping inactive records means that a withdrawn application and its
replacement can both appear. An earlier withdrawn application counts as an
alternative, not an additional building, when a valid archived BIN, the same
lot and the exact address identify one later non-withdrawn application within
365 days. Several withdrawn versions may share one successor. In the canonical
historical cohorts this rule excludes 34 earlier applications carrying 734
units from additive totals. The records keep the original filing date, the
successor's filing date and \texttt{refiling\_basis}; the archive gives no
withdrawal date, and none is inferred. Other inactive proposals remain
additive constituents.

Two committed decisions in \path{parent_opportunities_manual/output/historical_filing_roles.csv}
cover cases the rule cannot:

\begin{center}
\small
\begin{tabular}{llp{0.24\textwidth}rp{0.30\textwidth}}
\hline
Job & Address & Role & 23Q4 units & Basis \\
\hline
\texttt{B00775071} & 773 Neptune & Nonresidential filing & 189 &
23Q4 describes a one-story house of worship; the later DOB initial record
classifies it as a school with zero dwellings. \\
\texttt{M00536051} & 504 West 49th & Withdrawn alternative to \texttt{M00580473} & 158 &
Same lot and program as the permitted 158-unit 509 West 48th filing; HPD's
March 2022 environmental assessment describes one western housing building. \\
\hline
\end{tabular}
\end{center}

773 Neptune's 189 units remain as source evidence but add nothing to the
residential total. The separate residential job \texttt{B00646589} has another
BIN and lot. The West 49th filing, dated July 29, 2021, has a different BIN and
address from its successor, dated August 31, 2021, so the automatic rule does
not apply. The environmental assessment does not name either application, so
the link is a documented inference. The earlier date anchors the parent, and
the reviewed 21,924-square-foot allocation is unchanged. Together with the
automatic rule, 35 filings with 892 recorded units are now
superseded source filings.

\paragraph{BIN collisions.}
Distinct building identifiers among additive filings are part of the
composition-eligibility requirement for reweighting. Bedford Square has four
documented buildings, but the archive gives two of them the same BIN. The
saved official DOB initial records give all four distinct valid BINs, and the
complete filing set matches the reviewed production land allocation. A general
rule now admits a historical parent when both forms of corroboration hold; the
archived collision is kept and \texttt{reviewed\_distinct\_buildings} marks the
parent. Bedford is the only historical parent that qualifies. Its four
buildings, 877 units and 177,520.79 square feet remain in the weighting sample.

\paragraph{Exposure.}
The historical ownership and description screen also reads 23Q4. Attorney
General offering-plan searches cover every historical parent, but only plans
with dated submissions before the reference snapshot inform the
classification. These are 2026 searches of dated records, not an archived
2023 search. The supplement adds 652 exact job and address queries and one
cross-year query for the Livingston companion. A missing search is kept
distinct from a confirmed zero result. HPD 485-x registrations describe
post-policy projects only.

\paragraph{Resulting panels.}
The 23Q4 archive contains 2,432 new-building filings with at least six units
filed in 2019--2022: 2,166 active/completed-file records and 266 others. The
existing requirement of a positive-area parcel match before filing leaves 42
outside the historical linkage universe. Seven join parents anchored before
2019, and 36 remain in membership as nonadditive source filings. The other
2,347 additive filings, plus the 2023 Livingston companion inside its parent's
365-day window, are the canonical constituents.

\begin{center}
\small
\begin{tabular}{lrrr}
\hline
Panel & Parents & Additive constituents & Units \\
\hline
Historical, 25Q4 rule (before) & 1,811 & 1,904 & 129,335 \\
Historical, 23Q4 rule (current) & 2,173 & 2,348 & 141,632 \\
Post-policy (unchanged) & 907 & 1,005 & 61,319 \\
\hline
\end{tabular}
\end{center}

Of the new historical parents, 206 include an inactive additive constituent;
another 25 have only active additive constituents but an inactive earlier
application. The flags \texttt{historical\_all\_active} and
\texttt{historical\_all\_source\_active} record these two definitions. The
post-policy dates, memberships, exposure, composition eligibility and land
covariates were verified unchanged.

Livingston is now one parent with 144+105=249 units. The 23Q4 file records 372
Livingston, job \texttt{B00814498}, at 105 units, filed 120 days after 370.
Both archived filings share lot~16 and refer to the neighboring building under
a separate application. The earlier linkage used 25Q4 lots~22 and~16 and missed
the link; the later 99-unit design no longer enters the historical outcome.

\paragraph{Effect on the headline comparison.}
Historical constituent filings at exactly 99 units fall from six to
two. The reweighting balances log land area, residential FAR, built FAR
and borough across 595 historical and 310 post-period parents with at
least 50 units. When the source rule changed, the weighted historical share of
parents totaling exactly 99 units fell from 1.3537\% to 0.5261\%, against
11.9355\% in the post period. The rebuilt outputs now give
0.5245\% and 11.9355\%. This parent-total share includes one
historical parent whose smaller constituents sum to 99, so it differs from the
count of individual 99-unit filings. The historical effective sample size is
532.15, the maximum absolute balance error is
$1.48\times10^{-9}$, and 495 of 499 bootstrap draws calibrate,
with four failures recorded. These remain descriptive comparisons while
parcel boundaries are under review.

\paragraph{Consequences for the boundary queue.}
The source change altered historical membership, so the old parcel-review
queue was reconciled with the new sample. Of the 111 cases previously queued,
103 remain flagged; two of them now belong to one parent, giving 102 distinct
parents. The other eight left the weighting sample without being classified as
resolved, and the source correction cleared none of the old cases. At the time
of the change, 29 newly entering parents and 21 previously supported parents
were flagged as well, for 152 in total. After the later boundary decisions the
screen flags 145 of the 905 weighting parents
(116 historical and 29 post-period). The flags
identify coverage and allocation questions, not established measurement
errors. Documentary land reviews keep the filing sets they were made for. The
two Jamaica/Sutphin reviews cover separate filing sets that now form one
mechanically linked parent, so their separate allocations do not establish the
merged parent's boundary. West 48th is outside the weighting sample.

\paragraph{Remaining timing concern.}
The 23Q4 archive is a pre-adoption snapshot, not a record of each design as
first filed. Its counts reflect designs as of early 2024 and may include
amendments made after the initial 2019--2022 filing. Livingston's 144 and 105
are the counts recorded before adoption; the sources do not show that they
were the September 2022 counts. The rule removes post-adoption redesigns from
the historical distribution but does not recover original filings. A
filing-level history would be needed to separate the two.

\textit{Reproduction note.} Root \texttt{make data} rebuilds the panels and
\texttt{make reweighting} the weighting results in
\path{tasks/audits/audit_scale_shape_splitting/output}. The 23Q4 archive is a
checksum-pinned download in \path{tasks/fetch_dcp_housing_database}. Unit
selection is in \path{tasks/construct_parent_cohorts}, and reviewed roles are
in \path{tasks/parent_opportunities_manual}. The release comparison, the DOB
fallback listing, the check of the rebuilt panels against the dated pre-change
tables, and the queue reconciliation were one-time audits in
\path{tasks/audits/audit_hdb_mappluto_condo_recovery}, preserved at tag
\texttt{pre-cleanup-2026-09-22}. The earlier pilot's exhibits use the old
historical rule and are frozen in the logbook archive.
